\sscp{Subject indexing on verbs}{13-03-03}
\citet[258]{Blust1993} saw “the use of proclitic subject markers
on the verb” as a “morphosyntactic innovation” in support of PCEMP.
In chapters \ref{ch:CM}, \ref{ch:SB}, \ref{ch:AS}
it was shown that such markers are found in \tsc{Ambon-Seram}
languages, but not in most \tsc{Sula-Buru} languages,
and the forms that are not historical retentions from PMP are
not cognate with each other.
In the \tsc{Timor-Babar} subgroup (\crf{ch:TB}),
subject markers are found in all nodes,
but the forms that are not retentions are not cognate with each
other (discussed further below).
In the \tsc{Bima-Lembata} language of Dhao (\crf{ch:BL}),
one active motion verb ‘go’ is inflected with suffixes,
some with prefixes, and some not at all \citep{Grimes2010a}.
Subject prefixes or proclitics are also found widely west of
the Blust line, such as throughout much of Sulawesi (e.g.
ergative prefixes in South Sulawesi languages \citep{Friberg1991};
the complex subject prefixes in Muna \citep{vandenBerg1989},
and Busoa \citep[54ff]{vandenBerg2020}.
Several things can be observed regarding subject indexing on
verbs:

\begin{itemize}
\item	Their distribution is sporadic.
	\begin{itemize}
		\item	In languages that do have subject indexing,
					there may be many verbs that do not take it.
					In several subgroups verbs for ‘eat’
					and ‘drink’ are among those most likely to take subject indexing.
		\item	As noted in \srf{01-08},
					the forms that are not retentions from PMP are also not cognate
					with each other, both between Wallacean subgroups
					and even within the same subgroup.
					So, for example, \tsc{Ambon-Seram} forms
					are not cognate with the \tsc{Rote-Meto} forms.
					And within the \tsc{Timor-Babar} subgroup itself,
					the subject prefixes are not fully cognate with one another.
	\end{itemize}
\item	Some languages have no subject indexing. Others do.
\item	The presence of subject markers does not define any of
			our higher level subgroups (except perhaps \tsc{Ambon-Seram}),
			but it is a feature of some lower level groups.
\item	The forms that are cognate trace back to PMP
			and thus are retentions and not innovations.
			As such they are not useful for linguistic subgrouping.
\end{itemize}

\citet[59]{Blust2009} later concedes,
“If Wolff is correct (and the evidence suggests that he is),
then the use of prefixal/proclitic agreement markers on the verb
has arisen independently through multiple historical changes.
This obviously weakens it as subgrouping evidence.”

We also note in passing that affixing of core arguments is a
common feature of Papuan languages in the region.
\citet{Klamer2017} illustrates the many Undergoer prefixes found
on verbs in Undergoer-prominent \tsc{Alor-Pantar} languages which are typologically SOV.
Wersing is an \tsc{Alor-Pantar} language that has Undergoer prefixes
on many transitive verbs,
and also has subject (Actor) prefixing on some transitive
and intransitive verbs \citep{Jones2023}. 
The presence and forms of proclitic subject markers is of dubious
value for subgrouping purposes in the Wallacean region.
The difficulties in reconciling the diverse forms can be illustrated
by the \tsc{Timor-Babar} subgroup.
A sample of subject prefixes from \tsc{Timor-Babar} is given
in \trf{13-21}, which represents every branch of \tsc{Timor-Babar} (with two
examples from \tsc{Southwest Maluku}).

\begin{table}
\caption[Sample of \tsc{Timor-Babar} subject prefixes]{Sample of \tsc{Timor-Babar} subject prefixes\su{†}}\label{tab:13-21}
\begin{adjustbox}{width=\textwidth}
	\begin{threeparttable}\st{3pt}
	\begin{tabular}{ll|l|l|l|l@{ }l|l|ll|l}\lsptoprule
								&	Funai 			&	Tii		&	Tetun	&	Idate	&	\mc{2}{l|}{Tugun}				&	Luang	&	\mc{2}{l|}{SE Babar\su{Δ}}&	Selaru\\
								&	Helong\su{‡}&				&				&				&	\tsc{realis}&	\tsc{irr}	&				&						&								&	\\	\midrule
\tsc{1sg}				&	ku-					&	ʔa-		&	k-		&	u-		&	u-					&	mu-				&	u-		&	o-				&	iy-						&	ku-\\
\tsc{2sg}				&	mu-					&	ma-		&	m-		&	o-		&	o-					&	om-				&	mu-		&	mo-				&	my-						&	mu-\\
\tsc{3sg}				&	na-					&	i-/na-&	n-		&	na-		&	{\0}-				&	ma-				&	na-		&	le-				&	l-						&	i-/ki-\su{‖}\\
\tsc{1pl.incl}	&	ta-					&	ta-		&	{\0}-	&	ta-		&	it-					&	ka-				&	ta-		&	ke-				&	x-						&	ta-\\
\tsc{1pl.excl}	&	mi-					&	ma-		&	{\0}-	&	{\0}-	&	am-					&	am-				&	ma-		&	me-				&	m-						&	mi-\\
\tsc{2pl}				&	mi-					&	ma-		&	{\0}-	&	{\0}-	&	mi-					&	mar-			&	mi-		&	mi-				&	my-						&	mi-\\
\tsc{3pl}				&	na-					&	ra-		&	r-/n-	&	ra-		&	ra-					&	mar-			&	ra-		&	te-				&	t-						&	ra-\\
		\lspbottomrule
	\end{tabular}
		\begin{tablenotes}\tnsize
			\item[†]	{Most \tsc{Timor-Babar} languages have more than one set of prefixes.
								Which set a verb takes is determined by the phonotactics,
								and/or the semantics of the verb, and/or partially
								lexically determined. These factors vary from language to language.
								\trf{13-21} shows the most conservative prefixes for most languages.}
			\item[‡]	{Funai Helong subject prefixes are identical to Roi{\Q}is Amarasi,
								from which they may be borrowed.
								Semau Helong has lost productive prefixes,
								but retains fossilised \it{na-} on some verbs.}
			\item[Δ]	{SE Babar has undergone *n > \it{l},
								and *t > \it{k, x} thus partly explaining
								the \tsc{3sg} and \tsc{1pl.incl} prefixes.}
			\item[‖]	{Selaru \it{i-} is used for animate subjects
								and \it{ki-} for inanimate subjects.}
		\end{tablenotes}
	\end{threeparttable}
\end{adjustbox}
\end{table}

As can be seen from the diversity of forms in \trf{13-21},
it would be difficult to reconstruct subject prefixes to Proto-Timor-Babar,
except maybe \tsc{1sg} *ku-, and \tsc{2pl} *mi-.
Instead, individual branches have derived subject prefixes through
different grammaticalisation pathways. E.g.
some \tsc{2sg} prefixes are grammaticalised from *kahu (> **kau
> **o), while others are probably grammaticalised from \tsc{2pl} *kamu(-yu).
Given the difficulty in reconciling the forms in \tsc{Timor-Babar},
it is hard to see how we could reconstruct them for a higher
level subgroup such as PCMP.\footnote{
		A reviewer notes that many
		subject prefixes in Wallacea are consistent with a reconstruction
		\tsc{1sg} *ku-, \tsc{2sg} *mu-, \tsc{3sg} *na-,
		\tsc{1pl.incl} *ta-, \tsc{2pl} *mi-,
		\tsc{3pl} *da-, which are related to the PMP genitive enclitics
		and that the prefixes *ku-/*ta-/*da- often
		have prenasalised variants, which is compatible with a genitive-derived origin.
		The reviewer proposes that some of the discrepancies
		discussed in this section may be related to the emergence of
		a “second” generation of bound subject person-markers derived
		from free pronouns in subject position that fuse with the older subject prefixes.}
What is far more interesting with regard to subject-indexing
are the higher-level systems behind them,
such as widespread semantic alignment seen in the pronominal
prefixes/proclitics and suffixes/enclitics on verbs in many languages in the region.
There is a widespread distinction between Actor prefixes/proclitics
and Undergoer suffixes/enclitics, such as in \tsc{Ambon-Seram} languages (see \srf{07-04-03}).
This is reflected clearly in Split-S languages,
for example, in \tsc{Ambon-Seram}
and some \tsc{Timor-Babar} languages.
Yet Split-S languages of this sort are interspersed among languages
that have other typological or syntactic systems,
such as the nominative-accusative systems that have developed
in \tsc{Rote-Meto} languages (see Tables \ref{tab:13-17}, \ref{tab:13-18}, \ref{tab:13-19}).

